\documentclass{beamer}

\usepackage{aula}

\title{PyQGIS - plugins de interface}
\date{\today}

\begin{document}
    
\input{capa}

\section{QGIS Plugin Builder}

\begin{frame}{Plugin Builder: Passo a passo}
\begin{figure}
    \centering
    \includegraphics[width=0.5\linewidth]{pyqgis_plugins_figs/builder_menu}
\end{figure}
\begin{figure}
    \centering
    \includegraphics[width=0.7\linewidth]{pyqgis_plugins_figs/builder_install}
\end{figure}
\end{frame}

\begin{frame}{Plugin Builder: Passo a passo}
    \begin{figure}
        \centering
        \includegraphics[width=0.5\linewidth]{pyqgis_plugins_figs/builder_menu2}
    \end{figure}
    \begin{figure}
        \centering
        \includegraphics[width=0.7\linewidth]{pyqgis_plugins_figs/builder_wizard1}
    \end{figure}
\end{frame}

\begin{frame}{Plugin Builder: Passo a passo}
    \begin{columns}
        \begin{column}{0.48\linewidth}
            \begin{figure}
                \centering
                \includegraphics[width=\linewidth]{pyqgis_plugins_figs/builder_wizard2}
            \end{figure}
        \end{column}
    
        \begin{column}{0.48\linewidth}
            \begin{figure}
                \centering
                \includegraphics[width=\linewidth]{pyqgis_plugins_figs/builder_wizard3}
            \end{figure}    
        \end{column}
    \end{columns}
\end{frame}

\begin{frame}{Plugin Builder: Passo a passo}
    \begin{columns}
        \begin{column}{0.48\linewidth}
            \begin{figure}
                \centering
                \includegraphics[width=\linewidth]{pyqgis_plugins_figs/builder_wizard4}
            \end{figure}
        \end{column}
        
        \begin{column}{0.48\linewidth}
            \begin{figure}
                \centering
                \includegraphics[width=\linewidth]{pyqgis_plugins_figs/builder_wizard5}
            \end{figure}    
        \end{column}
    \end{columns}
\end{frame}

\begin{frame}[fragile]{Plugin Builder: Onde salvar?}
\begin{itemize}
    \item   Onde salvar?
    \item   LINUX
    \item   \url{/home/username/.local/share/QGIS/QGIS3/profiles/default/python/plugins/}
    \item   WINDOWS
    \item   \url{C:\\Users\\<user>\\AppData\\Roaming\\QGIS\\QGIS3\\profiles\\default\\python\\plugins\\}
\end{itemize}

    \begin{minted}{python}
from qgis.core import QgsApplication
print(QgsApplication.qgisSettingsDirPath())
    \end{minted}
\end{frame}

\begin{frame}[fragile]{Estrutura básica de um plugin QGIS}
    
    
    \begin{columns}
        \column[t]{0.48\textwidth}
        Principais componentes de um plugin PyQGIS:
        \begin{itemize}
            \item Classe principal do plugin
            \item Interface do QGIS (\texttt{iface})
            \item Arquivo de metadados
            \item Interface gráfica (Qt / .ui)
        \end{itemize}
        
        \column[t]{0.48\textwidth}
        
        \begin{figure}
            \centering
            \includegraphics[width=0.8\textwidth]{pyqgis_plugins_figs/plugin}
        \end{figure}
    \end{columns}
    
\end{frame}

\begin{frame}{Criação de Plugin: Estrutura básica do plugin}
    \begin{itemize}
        \item \_\_init\_\_.py = The starting point of the plugin. It has to have the classFactory() method and may have any other initialisation code.
        \item mainPlugin.py = The main working code of the plugin. Contains all the information about the actions of the plugin and the main code.
        \item resources.qrc = The .xml document created by Qt Designer. Contains relative paths to resources of the forms.
        \item resources.py = The translation of the .qrc file described above to Python.
        \item form.ui = The GUI created by Qt Designer.
        \item form.py = The translation of the form.ui described above to Python.
        \item metadata.txt = Contains general info, version, name and some other metadata used by plugins website and plugin infrastructure.
    \end{itemize}
    \url{https://docs.qgis.org/3.34/en/docs/pyqgis_developer_cookbook/plugins/plugins.html}
\end{frame}

\begin{frame}{Criação de Plugin: GitHub}
\textbf{Objetivo: Colocar a pasta .git e o código gerado pelo Plugin Builder na pasta do QGIS.}

\begin{center}
GitHub Desktop + QGIS + Plugin Builder
\end{center}
    \begin{itemize}
\item Criar pasta do plugin no Plugin Builder
\item Criar repositório GIT no GitHub
\item Baixar (Clone) o repositório para a máquina local
\item Copiar os arquivos de dentro da pasta do plugin para o repositório
\item Mover pasta do repositório para a pasta de plugins do QGIS
\item No GITHUB Desktop, adicionar repositório local e apontar para a pasta do plugin dentro da estrutura do QGIS
    \end{itemize}
\end{frame}

\section{Depuração remota no QGIS}

\begin{frame}{Depuração de plugins QGIS — Visão geral}
    
    Depurar plugins QGIS pode ser desafiador porque o código roda \textbf{dentro do QGIS}.
    
    \vspace{0.4cm}
    
    \textbf{Estratégia:}
    
    \begin{itemize}
        \item Usar o \textbf{VSCodium} como debugger
        \item Conectar ao processo do QGIS
        \item Inserir breakpoints no código do plugin
    \end{itemize}
    
    \vspace{0.4cm}
    
    \textbf{Ferramenta principal:}
    
    \begin{itemize}
        \item \texttt{debugpy} (debugger Python)
    \end{itemize}
    
\end{frame}

\begin{frame}[fragile]{Passo 1 — Inserir debugpy no plugin}
    
    Adicione no código do plugin (ex: no \texttt{initGui()} ou início do run):
    
    \begin{minted}{python}
import debugpy

debugpy.listen(("localhost", 5678))
print("Aguardando debugger...")

debugpy.wait_for_client()
print("Debugger conectado!")
    \end{minted}
    
    \vspace{0.3cm}
    
    \textbf{Efeito:}
    
    \begin{itemize}
        \item QGIS pausa execução
        \item Aguarda conexão do VSCodium
    \end{itemize}
    
\end{frame}

\begin{frame}[fragile]{Passo 2 — Configurar VSCodium}
    
    
\textbf{Ao iniciar a depuração, escolha "Remote attach".}
\begin{figure}
    \centering
    \includegraphics[width=0.7\linewidth]{pyqgis_figs/debug_remote_attach}
\end{figure}
    
    \vspace{0.3cm}
    
    \textbf{Depois:}
    
    \begin{itemize}
        \item Abrir o plugin no QGIS
        \item Executar a depuração.
    \end{itemize}
    
\end{frame}

\begin{frame}[fragile]{Passo 3 — Depuração na prática}
    
    Fluxo de uso:
    
    \begin{enumerate}
        \item Abrir QGIS
        \item Executar o plugin
        \item QGIS pausa em \texttt{wait\_for\_client()}
        \item Iniciar debug no VSCodium
        \item Continuar execução
    \end{enumerate}
    
    \vspace{0.4cm}
    
    \textbf{Agora você pode:}
    
    \begin{itemize}
        \item Colocar breakpoints
        \item Inspecionar variáveis
        \item Executar passo a passo
    \end{itemize}
    
    \vspace{0.4cm}
    
    \textbf{Dica:}
    
    Remova os comandos antes de publicar o plugin.
    
\end{frame}


\end{document}
